\documentclass[a4paper,11pt]{article}
\usepackage[utf8]{inputenc}
\usepackage[T1]{fontenc}
\usepackage[english]{babel}
\usepackage{geometry}
\geometry{left=2.5cm,right=2.5cm,top=2cm,bottom=2.5cm}

\usepackage{lmodern}
\usepackage{xcolor}
\definecolor{titleblue}{RGB}{0,51,140}
\usepackage{hyperref}

\pagestyle{empty}

\begin{document}

\begin{flushleft}
    \textbf{Ismail AISSAOUI} \\
    State Engineer in Artificial Intelligence \\
    \href{mailto:ismail.aissaoui.pro@gmail.com}{ismail.aissaoui.pro@gmail.com} \\
    +213 660 70 77 96 \\
    \href{https://ismail-aissaoui.vercel.app}{ismail-aissaoui.vercel.app}
\end{flushleft}

\begin{flushright}
    To P.M. Congedo, E. Denimal Goy, and Olivier Le Maître \\
    \textbf{Inria Platon Team / CMAP} \\
    École Polytechnique \\
    \medskip
    May 12, 2026
\end{flushright}

\bigskip

\noindent \textbf{Subject: Application for PhD in Multi-Fidelity Scientific Machine Learning for Digital Twins}

\bigskip

Dear Members of the Selection Committee,

As a State Engineer in Artificial Intelligence from ESI-SBA, specializing in the intersection of physics and deep learning, I am highly motivated to apply for the PhD position on "Multi-Fidelity Scientific Machine Learning with Heterogeneous Inputs" within the Platon project-team. Your research focus on bridging information across different fidelity levels perfectly aligns with my professional goal of developing robust and scalable Digital Twins for critical systems.

During my Master’s thesis at Sonatrach, I designed a **Physics-Informed Digital Twin** for gas turbine monitoring. The core challenge of my work was to build a generative surrogate model capable of accurately predicting complex thermodynamic states while maintaining real-time performance. By leveraging **Mamba-SSM** and **xLSTM** architectures and implementing custom **PINN** loss functions, I achieved a balance between physical consistency and a latency of \textbf{0.04ms}. This experience taught me that the biggest hurdle in industrial Digital Twins is not just data volume, but the intelligent fusion of multi-fidelity signals—ranging from coarse sensors to high-fidelity physical simulations.

The proposed PhD topic resonates with my background because it addresses the fundamental problem of **heterogeneity**. In my current work, I have dealt with mismatched data structures between low-fidelity sensor streams and high-fidelity physics-based solvers. I am particularly excited about the axes involving **latent-variable models (VAEs)** and **Operator Learning (DeepONet, FNO)** to align these heterogeneous spaces. My experience with hardware-software co-design and C++ JIT optimization ensures that I can not only design these theoretical frameworks but also implement them into the efficient computational pipelines required by the **MediTwin** project.

Joining the CMAP at École Polytechnique and the Inria Platon team would provide the ideal environment for me to master the mathematical rigors of uncertainty quantification and information fusion. I am a highly autonomous researcher with a strong background in PyTorch and numerical optimization, and I am eager to contribute to the development of patient-specific digital twins that are both reliable and computationally tractable.

Thank you for your time and consideration. I look forward to the possibility of discussing my research background and how my skills in physics-aware AI can contribute to the success of this project.

\bigskip

Sincerely,

\bigskip

\textbf{Ismail AISSAOUI}

\end{document}
